\subsection*{4.2 A reduction from $r$-Bounded Set Cover to Multi-MCSP}
In order to show that Multi-MCSP is NP-hard, we give a probabalistic polynomial-time many-one reduction with one-sided error from a constant approximation of $r$-Bounded Set Cover to Multi-MCSP.
In fact, for convenience, our reduction will be from the optimization version of approximating $r$-Bounded Set Cover to the optimization version of Multi-MCSP (computing $\CCirc$), but it will be easy to see that our reduction can be converted into the desired reduction for the corresponding decision problems.
Let $r$ be a large enough constant, so that say 10-approximating $r$-Bounded Set Cover is NP-hard. Given an instance $(1^n,\mathcal S)$ of this problem, the reduction proceeds as follows. Let $m=O(n^3)$ be the least power of two greater than $n^3$. Let $T$ be a uniformly random truth table of length $m$. I.e., $T$ is a binary string in $\{0,1\}^m$, representing a function from $\{0,1\}^{\log m}$ to $\{0,1\}$. Compute the truth table of
\[g\stackrel{\mathrm{def}}=\mathop{\bullet}_{S\in\mathcal S}\operatorname{Eval}\text{-}\operatorname{DNF}_{T\langle S^m\rangle}.\]
Let $k$ be the number of distinct components of $g$ that are not functions computed by an input gate, that is,
\[k\stackrel{\mathrm{def}}=|\{g_i:g_i\text{ is a component of }g\text{ and }g_i\ne x_j\text{ for all }j\in[\log m]\}|.\]
The reduction then outputs\footnote{For the decision problems, we can determine if there is a 10-approximate set cover of size $\ell$ by outputting $\operatorname{Multi}\text{-}\operatorname{MCSP}(T\bullet g,k+\ell)$, which computes whether $\Delta\leq\ell$.}
\[\Delta\stackrel{\mathrm{def}}=\CCirc(T\bullet g)-k.\]
First, we argue that this procedure runs in polynomial time. The only two steps that may raise concern are whether we can compute the truth table of $g$ efficiently and whether we can compute $k$ efficiently.
To show that the truth table of $g$ can be computed efficiently, it suffices to show that for each $S\in\mathcal S$, the truth table of $\operatorname{Eval}\text{-}\operatorname{DNF}_{T\langle S^m\rangle}$ can be computed in time polynomial in $n$. Computing $T\langle S^m\rangle$ can be done in time $O(m+|T|)=O(m)$ and outputs a truth table of length $m$. The canonical DNF of the truth table $T\langle S^m\rangle$ can be computed in time polynomial in $|T\langle S^m\rangle|=m$, and the resulting DNF has $\log m$ inputs and size at most $O(m\log m)$. Finally, computing the Evaluation Function of a circuit with $\log m$ inputs and $O(m\log m)$ gates can be done in time $O(m^3)$ by just evaluating the circuit on every input. Hence, putting these all together, computing the truth table of $\operatorname{Eval}\text{-}\operatorname{DNF}_{T\langle S^m\rangle}$ can be done in time polynomial in $m=O(n^3)$.
To see that we can compute $k$ efficiently, realize that we have already computed the full truth table of $g$ and that removing any components computed by one of the $\log m$ input wires along with any duplicate components takes time at most quadratic in the length of the truth table of $g$.
Now, we will argue for the correctness of the reduction. We will show that
\[\underbrace{\operatorname{cover}([n],\mathcal S)/4-4}_{\substack{\text{with high probability}\\\text{using Lemma16}}}\ \leq\ \Delta\ \leq\ \underbrace{\operatorname{cover}([n],\mathcal S)}_{\substack{\text{unconditionally}\\\text{using Lemma15}}},\]
and thus, with high probability $\Delta$ computes a 10-approximation of $r$-Bounded Set Cover when $n$ is sufficiently large.\footnote{Since $\operatorname{cover}([n],\mathcal S)\geq n/r$ (using that the sets in $\mathcal S$ have cardinality at most $r$), we will actually get that, for each $\epsilon>0$, $\Delta$ gives a $4+\epsilon$ approximation with high probability when $n$ is sufficiently large.} Moreover, this computation has one-sided error since the upper bound holds unconditionally.
Thus, after proving Lemmas 15 and 16, we will have shown that there is a randomized polynomial-time many-one reduction with one-sided error from 10-approximating $r$-Bounded Set Cover to Multi-MCSP.
Before proving Lemmas 15 and 16, we make the following observation about computing $g$.
\setcounter{nttheorem}{13}
\begin{ntproposition}[Author Proposition14]
If a circuit $C$ computes $g$, then there are $k$ distinct gates in $C$ that compute components of $g$. Moreover, $\CCirc(g)=k$.
\end{ntproposition}
\begin{proof}
We begin by proving the first statement, which also implies the lower bound $\CCirc(g)\geq k$. Suppose $C$ is a circuit that computes $g$. Then every distinct component of $g$ has a (necessarily distinct) input wire or gate from $C$ that computes that component. Therefore, since $g$ has $k$ distinct components that are not computed by an input wire, $C$ must have at least $k$ distinct gates computing components of $g$.
Next, we sketch the proof of the upper bound $\CCirc(g)\leq k$. Let $C$ be the circuit built as follows. For each $S\in\mathcal S$, iterate through the gates $g$ in $\operatorname{DNF}_{T\langle S^m\rangle}$ in topological order. Let $\diamond\in\{\land,\lor,\neg\}$ be the gate type of $g$. If $g$ computes a function that is already computed by $C$, then ignore it. Otherwise, add a $\diamond$ gate to $C$ that takes as input(s) those gate(s) in $C$ that compute the function(s) which are fed as inputs to $g$ in $\operatorname{DNF}_{T\langle S^m\rangle}$ (we are guaranteed to find such gates in $C$ since we are iterating in topological order).
By construction, $C$ computes $g$. (Recall that $g=\mathop{\bullet}_{S\in\mathcal S}\operatorname{Eval}\text{-}\operatorname{DNF}_{T\langle S^m\rangle}$, and our construction ensures $C$ computes every function computed by a gate in $\operatorname{DNF}_{T\langle S^m\rangle}$ for any $S\in\mathcal S$.) Moreover, our construction maintains that every gate in $C$ computes a component of $g$ that is not computed by any other gate or input wire. Thus, since $g$ has at most $k$ unique components not computed by input wires, $C$ has at most $k$ gates.
\end{proof}
One consequence of Proposition14 is that $\Delta=\CCirc(T\bullet g)-\CCirc(g)$. With this fact, we can prove our two main lemmas.
\begin{ntlemma}[Author Lemma15]
If $\operatorname{cover}([n],\mathcal S)=\ell$, then $\Delta\leq\ell$.
\end{ntlemma}
\begin{proof}
Let $S_1,\ldots,S_\ell$ be a cover of $[n]$ using sets from $\mathcal S$. Then, by construction, we have that $T=T\langle S_1^m\rangle\lor\cdots\lor T\langle S_\ell^m\rangle$. Since $T\langle S_1^m\rangle,\ldots,T\langle S_\ell^m\rangle$ are components of $g$, this implies that
\[\Delta=\CCirc(T\bullet g)-\CCirc(g)\leq\ell\]
as desired.
\end{proof}
